\section{The exact remaining branch datum for the Gaussian scalar $E_8$ source}
\label{sec:e8-source-gate}

\subsection*{HD-E8-01. Local determinants in $\mu_4W(E_8)$}
Let $W=W(E_8)$ in its rational reflection representation, endowed
with the positive definite root form $H$. Use the simple-root basis
whose Cartan matrix has diagonal entries two and off-diagonal entries
minus one on the edges
\[
 (1,3),\ (3,4),\ (2,4),\ (4,5),\ (5,6),\ (6,7),\ (7,8).
\]
The determinant of $H$ is one. Let $s_j=I-e_jH_{j,*}$ and put
$G=\langle W,iI\rangle\subset U(H,\mathbf Q(i))$.
Since $-I\in W$, this group has two cosets $W$ and $iW$,
and order $2\cdot696729600$.

For $a\in G$ let $U_a=\operatorname{im}(1-a)$ and define, as in
D-G31-01,
\[
 k(a)=\dim U_a,\qquad
 q_+(a)=\frac{\det(H|_{U_a})}{|\det(1-a|_{U_a})|}
       \quad\bmod\operatorname{Nm}_{\mathbf Q(i)/\mathbf Q}\mathbf Q(i)^*.
\]
Every class in $W$ has norm $q_+$. Every class in $iW$ has moving
rank at least four. Its unique rank-four class is $i$ times Carter's
$2D_4(a_1)$, has age two, and has norm $q_+$. The complete list of
nonnorm classes in $iW$ is the following. The first column uses
the CHEVIE author's numbering, reproduced by the accompanying script.
All six classes have moving rank eight and local norm class three.
\[
\begin{array}{r|l|r|r|r}
\text{index}&\text{class of }w\text{ in }a=iw&
 \det(1-a)&q_+(a)&\operatorname{age}(a)-4\\\hline
55&E_6+A_2&-3&3&1\\
56&E_8(a_7)&-3&3&-1\\
57&E_6&6&6&0\\
58&E_7(a_2)+A_1&6&6&0\\
101&E_6+A_1&6i&6&1/2\\
102&E_7(a_2)&-6i&6&-1/2
\end{array}
\]
Here $6$ and $3$ have the same Gaussian norm class since $2$ is a norm.

\paragraph{Exact computation and its external input.}
The complete class list is the published Weyl-group classification
of Carter, \emph{Conjugacy classes in the Weyl group}, Compositio
Math.~25 (1972), 1--59, Proposition~21/Table~3, pp.~22--23, and
Table~11, pp.~52--54. We use the CHEVIE author's explicit complete
class representatives, documented in the GAP3 manual, section~87.1,
\emph{ChevieClassInfo}, and retained as mathematical words in
\texttt{research/e8\_class\_exploration.json}. Their retrieval URL
and SHA-256 occur in that file. Completeness of this known class
classification is an input, not a new exhaustive enumeration of
the $696729600$ group elements.

\texttt{scripts/examine\_e8\_classes.py} reconstructs each representative
as a product of the displayed integral simple reflections. It checks
all 112 matrix orders, orthogonality, and characteristic polynomials;
the supplied class sizes sum to the group order.
\texttt{scripts/e8\_gaussian\_local.py} computes a basis of every
moving space by exact Gaussian elimination, restricts both the
Hermitian form and the operator, and takes their determinants.
Its norm test factors rational numbers and tests the parity of
valuations at primes congruent to three modulo four. Ages are
computed from exact cyclotomic factors of the real characteristic
polynomial, with every exponent shifted by $1/4$ modulo one.
\texttt{scripts/e8\_real\_local.py} separately checks all real classes.
The complete rational outputs, including spectra, are retained.
This proves the displayed finite arithmetic assertions from the
specified complete representative list.

\subsection*{HD-E8-02. A necessary twelve-branch pattern}
Suppose the natural representation of $G$ supplies a dominant
nonsplit Weil-sixfold period map. Then the base is genus zero,
its signed Hermitian discriminant has Gaussian norm class three,
and its inertia tuple consists of exactly
\[
 \boxed{\text{ten real reflections, one rank-four element of }iW,
        \text{ and one rank-eight element }iw\text{ of class57 or58}.}
\]
The rank-four element is an involution, and the rank-eight element
has order twelve. In particular its branch orders are $(2^{11},12)$.
This is a necessary pattern. It does not assert product-one
generation, actual rank nine, or existence of detecting cycles.

\begin{proof}
For a base of genus $\gamma$, the dimension equation is
$6=(2\gamma-2)8+D$, where $D$ is the sum of local moving ranks.
For $\gamma\ge2$ this is impossible. At genus one $D=6$ and
the Hurwitz dimension is at most six. Dominance therefore requires
genus zero, $D=22$, and at least twelve branch points.

The signed determinant formula D-DET-01 applies. Since the space
is balanced, its determinant representative is positive; as in
D-G31-02 its norm class is the product of the positive local factors
$q_+$. Thus nonsplitness requires one of the six exceptional entries
in HD-E8-01, of moving rank eight. A product-one tuple has an even
number of entries in the nonreal coset $iW$, and a generating tuple
has at least one. Thus another nonreal entry has moving rank at
least four. If there are $r$ entries, their total moving rank is at
least $8+4+(r-2)=r+10$. With $D=22$ and $r\ge12$, equality
is forced: $r=12$, the second outside entry has rank four, and
every other entry is a reflection. Such reflections belong to $W$.

Each real reflection has age $1/2$, and the rank-four outside
involution has age two. Balance requires total age $11$, hence the
rank-eight element must have age four. The table leaves precisely
classes57 and58. Its local norm class is three, all other entries
have norm local factors, and the stated discriminant follows.
The spectra in the certificate give order twelve in both cases.
\end{proof}

\paragraph{Source URLs.}
\url{https://www.numdam.org/item/CM_1972__25_1_1_0.pdf};
\url{https://webusers.imj-prg.fr/~jean.michel/gap3/htm/chap087.htm};
\url{https://raw.githubusercontent.com/jmichel7/gap3-jm/master/pkg/chevie/tbl/weyle8.g}.
No entire third-party library file is redistributed.
